\documentclass[10pt,a4paper]{amsart}
\usepackage[margin=19mm]{geometry}
\usepackage{amsmath,amssymb,mathtools}
\usepackage[scr=rsfs,cal=euler]{mathalfa}
\usepackage{xfrac}
\usepackage[unicode,hidelinks,bookmarks=false]{hyperref}
\newcommand{\ie}{i.e.\ }
\newcommand{\cf}{cf.\ }
\newcommand{\resp}{respectively\ }
\newcommand{\Subsection}{\subsection}
\providecommand{\nobreakdash}{\nobreak\hbox{-}}
\setlength{\emergencystretch}{2em}
\multlinegap0pt
\usepackage{bm}
\usepackage{bbm}
\usepackage{dsfont}
\usepackage{xspace}
\usepackage{cancel}
\usepackage{mathtools}
\usepackage{amsthm}
\makeatletter
\@ifundefined{defin}{\newtheorem{defin}{Defin}}{}
\@ifundefined{lem}{\newtheorem{lem}[defin]{Lemma}}{}
\@ifundefined{prop}{\newtheorem{prop}[defin]{Proposition}}{}
\makeatother
\newcommand{\Alt}{\ensuremath{\operatorname{Alt}}}%
\newcommand{\Sym}{\ensuremath{\operatorname{Sym}}}%
\newcommand{\Z}{\ensuremath{\mathbb{Z}}}%
\newcommand{\alt}{\ensuremath{\operatorname{Alt}}}%
\newcommand{\sub}{\ensuremath{\operatorname{Sub}}}%
\newcommand{\vol}{\ensuremath{\operatorname{vol}}}%
\title[On the growth of actions of free products]{On the growth of actions of free products}
\author{Adrien Le Boudec, Nicol\'as Matte Bon, Ville Salo}
\date{Source: arXiv:2401.06886v2 (CC BY 4.0)}
\begin{document}
\maketitle

\noindent\emph{Source and licence.} On the growth of actions of free products. Version arXiv:2401.06886v2 (CC BY 4.0), distributed under the Creative Commons Attribution 4.0 International licence, as recorded in the arXiv licence metadata for this exact version and shown on the abstract page https://arxiv.org/abs/2401.06886v2. Full rights notice: licenses/arxiv-2401.06886-CC-BY-4.0.txt.\par\smallskip

\noindent\textit{Editorial scope.} Counted unit: the Prop whose source label is \texttt{p-Houghton} in the article (source0.tex lines 711--722, source label \texttt{p-Houghton}), delivered together with the article's own complete original proof of it (source0.tex lines 724--744, one \texttt{proof} environment of 21 lines). No mathematics is written, repaired or re-derived: statement and proof are the article's own text line for line. Required context retained uncounted: lem:SmallGrowthConfined (lines 655--657); lem:FiniteIndex (lines 675--678); p-confined-partially-finitary (lines 696--698). Every definition, display and label of those lines is retained unchanged. Omitted from this package and not redistributed: 1--654, 658--674, 679--695, 699--710, 723--723, 745--837. Nothing inside a retained excerpt is omitted or altered: every retained body is delivered exactly as the article's lines are written, so no omission receipt of any kind accompanies this package. Citations: the retained excerpts carry no citation command at all, so no bibliography entry of the article is transcribed and no reference is redistributed. Reference adaptation: none was needed and none was made; every reference inside the delivered unit resolves inside it, so no unresolved reference or citation is produced. Printed slips kept literal: no printed word, symbol, hypothesis or number of the counted statement or of its proof was altered. Layout and class: the article's own class is \texttt{amsart} with packages including amssymb, amsfonts, amsthm, amsmath, mathrsfs, xspace, hyperref, graphicx, babel, inputenc, fontenc, geometry, csquotes, color; the wrapper is the lane's pinned \texttt{amsart} editorial wrapper at 10pt on A4, which restates the theorem-like environments the retained text needs and adds the article's own macro definitions that the retained text needs (\texttt{\textbackslash Alt}, \texttt{\textbackslash Sym}, \texttt{\textbackslash Z}, \texttt{\textbackslash alt}, \texttt{\textbackslash sub}, \texttt{\textbackslash vol}), with extra wrapper packages (bm, bbm, dsfont, xspace, cancel, mathtools). The delivered numbering is the unit's own. Assets: no retained span contains a figure environment, a graphics-inclusion command, a PDF-page inclusion, a table or a TikZ picture; no image file is copied into this package, the package declares no asset dependency and no image file is redistributed with it. Licence: version arXiv:2401.06886v2 of this article is distributed under the Creative Commons Attribution 4.0 International licence, as recorded in the arXiv licence metadata for this exact version and shown on the abstract page https://arxiv.org/abs/2401.06886v2. A full rights notice is delivered at licenses/arxiv-2401.06886-CC-BY-4.0.txt. Characters: every retained excerpt is pure ASCII, so the delivered engine, XeLaTeX with its default Latin Modern text font, reports no missing character.
\begin{lem}\label{lem:SmallGrowthConfined}
Let $G$ be a finitely generated group, and $X$ a $G$-set. Let \[ \mathcal{S}(X) = \overline{ \left\lbrace G_x  \, : \, x \in X \right\rbrace } \subseteq \sub(G). \] Then for every $H \in \mathcal{S}(X)$, we have $\vol_{G, G/H}(n) \preccurlyeq \vol_{G, X}(n)$. In particular if some $G_x$ is not confined then $\vol_{G,X}(n) \simeq \vol_G(n)$. 
\end{lem}

\begin{lem}
	\label{lem:FiniteIndex}
	Let $G$ be a finitely-generated group, and let $H \leq K$ be subgroups of $G$. Then $\vol_{G, G/K} \le \vol_{G, G/H}$, and if in addition $H$ has finite index in $K$, then $\vol_{G, G/H} \simeq \vol_{G, G/K}$.
\end{lem}

\begin{prop}\label{p-confined-partially-finitary}
Let $\Omega$ be a countable set, and let $G \leq \Sym(\Omega)$ be any subgroup containing $\Alt_{\mathrm{f}}(\Omega)$.  Then a subgroup $H\le G$ is confined if and only if $\Omega_{H, \mathrm{f}}$ is finite and there exists an $H$-invariant partition $\Omega_{H,\infty} = \Omega_1 \cup \cdots \cup \Omega_k$ such that $H$ contains $\Alt_{\mathrm{f}}(\Omega_1) \times \cdots \times \Alt_{\mathrm{f}}(\Omega_k)$.
\end{prop}

\begin{prop} \label{p-Houghton}
For $r \ge 2$, let $X$ be a faithful $H_r$-set such that $\vol_{G, X}(n) \not \succeq n^2$. Let $X=\sqcup_{i \in I} X_i$ be the decomposition of $X$ into  $H_r$-orbits, and let $J$ be the subset of $i\in I$ such that the $H_r$-action on $X_i$ is  faithful.
Then:
\begin{enumerate}
\item \label{i-A-trivial} for $i\notin J$, the $A$-action on $X_i$ is trivial;
\item \label{i-A-non-trivial} there exist finite  $H_r$-sets $(Y_i)_{i\in J}$ of bounded cardinality (with respect to $i$) such that for every $i\in J$, the $H_r$-set $X_i$ is isomorphic to the product $H_r$-set $\Omega\times Y_i$.


\end{enumerate}
In particular, if $Z=\bigcup_{i\notin J} X_i$ is the set of $A$-fixed points, there exists a finite index subgroup $K$ of $H_r$ such that every $K$-orbit in $X\setminus Z$ is isomorphic to $\Omega$ as a $K$-set.
%for every $i$ either the action of $G$ on $X_i$ is not faithful (equivalently, the group $A$ acts trivially), or  it is conjugate to the diagonal action on $\Omega \times Y$, where $Y$ is a finite $G$-set whose cardinality is uniformly bounded (that is, the upper bound that does not depend on the choice of the orbit $X_i$). 
\end{prop}

\begin{proof}

Set $G=H_r$. First note that item \eqref{i-A-trivial} follows from the fact that $A$ is contained in every non-trivial normal subgroup of $G$. We will show \eqref{i-A-non-trivial}, which is the content of the statement.
We fix a word metric $\lVert \cdot \rVert$ on $G$, associated to some symmetric generating set. Let $X$ be  a $G$-set as in the statement and $K$ be the stabiliser of a point $x_0\in X_i$, with $I\in J$. The group $G$ has exponential growth, so by Lemma~\ref{lem:SmallGrowthConfined}, $K$ must be confined. Thus, we can apply Proposition \ref{p-confined-partially-finitary}.

Let us first show that $|\Omega_{K, \mathrm{f}}|\le 1$. Suppose by contradiction that this is not the case. Let $K_0\le K$ be the pointwise stabilizer of the finite set $\Omega_{K, \mathrm{f}}$, which has finite index in $K$. Then $\vol_{G, G/K_0}(n)\simeq \vol_{G, X_i}(n)$ by Lemma~\ref{lem:FiniteIndex}.  Choose two distinct points $x_1, x_2\in \Omega_{K, \mathrm{f}}$ and let $K_1 \ge K_0$ be their pointwise stabiliser. Since $\vol_{G, X_i}(n) \succcurlyeq \vol_{G, G/K_1}(n)$, a contradiction will follow if we show that $\vol_{G, G/K_1}(n)\succcurlyeq n^2$. Since the action of $G$ on $\Omega$ is highly transitive (that is, transitive on $n$-tuples of distinct points), upon replacing $K_1$ by some $G$-conjugate we suppose that $x_i$ is the point at position 1 on $\ell_i$ for $i=1, 2$ (with the convention that the common origin is at position 0).

Now for $i=1, 2$, let  $\sigma_i$ be the transposition that swaps the first and second positions of $\ell_i$. Let also $t_i\in G$ be any element such that $t_i|_{\ell_i}$ shifts $\ell_i$ inside itself by 1. For example, these can be taken to be $t, t^{-1}$ where $t$ is the shift along a $\Z$-isomorphic ray coming from juxtaposing the rays $\ell_1, \ell_2$. For $i=1, 2$ and $n\ge 1$ consider the element
 \[\gamma_{i, n}=(t_i^n\sigma_i t_i^{-n})(t_i^{n-1} \sigma_i t^{-n+1}) \cdots \sigma_i=t_i^n\sigma_i (t_i^{-1}\sigma_i)^n.\]
Note that $\lVert\gamma_{i, n}\rVert \le Cn$, with $C=\max\{\lVert \sigma_i \rVert, \lVert t_i\rVert \;|\; i=1,2\}$.  On the other hand each $\gamma_{i, n}$ is a permutation with finite support contained in $\ell_i$, and $\gamma_{i, n}(x_i)$ is the point at position $n+1$ on $\ell_i$. Hence applying products of the form $\gamma_{1, m_1} \gamma_{2, m_2}$ to the pair $(x_1, x_2)$, with $0\le m_1, m_2\le n$ shows that $\vol_{G, G/K_1}(n)\succcurlyeq n^2$. This is a contradiction and proves the claim that $|\Omega_{K, \mathrm{f}}|\le 1$.

We now consider the partition $\Omega_{K, \infty}=\Omega_1\cup \cdots \cup\Omega_k$ and claim that $k=1$.  Suppose that $k\ge 2$. By a similar reasoning as above (passing first to a finite index subgroup of $K$, and then to an overgroup), we have $\vol_{G, X}(n)\succcurlyeq \vol_{G, G/K_1}(n)$, where $K_1$ is the subgroup of $G$ of elements that preserve both $\Omega_1$ and $\Omega_2$, so a contradiction will follow if we show that $\vol_{G, G/K_1}(n)\succcurlyeq n^2$. Fix $n$ large enough. Again by high transitivity  and using the fact that $\Omega_1, \Omega_2$ are both infinite, we can find $g\in G$ such that the point at position 1 on $\ell_1$ is in $g(\Omega_1)$, and all points on $\ell_1$ at positions between 2 and $n$ are in $g(\Omega_2)$, while the symmetric statement holds for $\ell_2$. Now applying to $g(\Omega_1)$ and $g(\Omega_2)$ the same the elements $h=\gamma_{1, m_1}\gamma_{2, m_2}$ defined above with $1\le m_i\le n$, we obtain $n^2$ different pairs $(hg(\Omega_1), hg(\Omega_2))$ with $\lVert h\lVert \le 2Cn$. This shows that the ball of radius $2Cn$ around $gK_1$ has cardinality at least $n^2$. Hence $k=1$. 

It follows that $K$ contains $\alt(\Omega_{K, \infty})$. This also shows that $\Omega_{K, \mathrm{f}}\neq \varnothing$, since otherwise $K$ would contain the normal subgroup $A$, contradicting that the action on $X_I$ is faithful. Hence $K\cap A$ contains a point stabiliser for the action of $A$ on $\Omega$, and hence is equal to it, as the latter is a maximal subgroup of $A$. Since $K$ was an arbitrary point stabiliser, we deduce that the action of $A$ on $X_i$ is conjugate on each of its orbits to the standard action of $A$ on $\Omega$. For each $A$-orbit $\mathcal{O}\subset X_i$ we denote by $j_\mathcal{O}\colon \mathcal{O}\to \Omega$ the unique $A$-equivariant bijection (its uniqueness follows from the fact that the only $A$-equivariant permutation of $\Omega$ is the identity). 

Let $Y_i=X_i/A$, the space of $A$-orbits, on which $G$ acts since $A$ is normal. For a point $x\in A$, we denote by $\mathcal{O}(x)$ its $A$-orbit. Consider the map $\iota\colon X_i \to \Omega \times Y_i$ given by $\iota(x)=(j_{\mathcal{O}(x)}(x), \mathcal{O}(x))$. This map is equivariant (this is obvious for the second component; for the first, it follows from the observation that the map $g^{-1}\circ j_{\mathcal{O}(gx)} \circ g$ is an $A$-equivariant bijection from $\mathcal{O}(x)$ to $\Omega$, and thus it is equal to $j_{\mathcal{O}(x)}$). The map $\iota$ is also clearly injective. To check that it is surjective it is enough to observe that the diagonal $G$-action on $\Omega \times Y_i$ is transitive: this is true because the $G$-action on $Y_i$ is transitive, and the group $A$ acts trivially on $Y_i$ and transitively on each fiber $\Omega\times \{y\}$.

To conclude we need to argue that  $Y_i$ must have uniformly bounded cardinality when $i\in J$ varies. Fix $n$ and suppose that the diameter of $Y_i$ is at least $n$. Then we can find a point $y\in Y_i$ and elements $s_1, \ldots, s_n$ in the generating set of $G$ such that, if we set $g_m=s_m\cdots s_1$, then the points $y, g_1y,  \ldots, g_ny$ are pairwise distinct (the point $y$ and  sequence $(s_i)$ can be found by looking at any path in the Schreier graph of $Y_i$ of length at least $n$). Next consider again the elements $\gamma_{i, n}$ above. Note that for every $n$, the element $\delta_n=\gamma_{1, n} \gamma_{2, n}$ belongs to $A$ and thus acts trivially on $Y_i$, and coincides with $\gamma_{1, n}$ in restriction to $\ell_1\subset \Omega$. Let $x\in \ell_1$ be at position 1, and consider the point $(x, y)\in \Omega \times Y_i$. Then the $n^2$ points $g_{m_2} \delta_{m_1}(x, y)=(g_{m_2}\gamma_{1,m_1}x, g_{m_2}y)$ for $m_1, m_2\le n$ are pairwise distinct, showing that there are at least $n^2$ in a ball of radius $O(n^2)$. This cannot be true for $n$ arbitrarily large, showing that the cardinality of $Y_i$ must be uniformly bounded. 

Finally, since the group $G$ is finitely generated, we may choose a finite index subgroup $K$ which acts trivially on $Y_i$ for every $i$, showing the last sentence in the statement. \qedhere
\end{proof}

\end{document}
